\documentclass[10pt]{article}
\usepackage[margin=20mm]{geometry}
\usepackage{amsmath,amssymb,mathtools,fontspec,unicode-math,xurl,hyperref}
\IfFontExistsTF{DejaVu Serif}{\setmainfont{DejaVu Serif}}{\setmainfont{Latin Modern Roman}}
\setmathfont{latinmodern-math.otf}
\setlength{\parindent}{0pt}\setlength{\parskip}{6pt}\sloppy
\title{Canonical Metrics, Orbifold Bundles, and Deformation Arguments}
\author{AI-assisted technical review}
\date{}
\begin{document}
\maketitle
Abstract: The frozen author version arXiv:2304.02582v1 is reviewed. Four Major findings concern the metric construction and finite descent in Theorem 2.4, the inclusion of zero-Lee metrics in the printed definition, the principal-bundle image criterion, and the cohomological step in negative-type deformation stability. Eight Minor corrections concern effective orbifolds, coefficients, normalization, deformation coordinates, curvature, and examples. Explicit witnesses and repair boundaries are given. The principal conclusions for strict Vaisman manifolds are not refuted by these checks. A separate conditional issue about normal Heisenberg lattices remains a limitation of the group classification assessment. The scientific decision is pending.\par
\section*{Source identity and mathematical scope}
Frozen identifier: M25-0260. Author: Nicolina Istrati. Article: Vaisman manifolds with vanishing first Chern class. The frozen DOI field is\par
https://doi.org/10.4171/jems/1633.\par
The actual source is the 26-page author manuscript arXiv:2304.02582v1, submitted on 5 April 2023 at 16:51:36 UTC. Its 531 652 bytes have SHA256\par
`be2abde0bac13aadcd414829ef38a62b58011486dda3a3f994d21a14a8cb492a`.\par
The publisher records submission on 22 June 2023, acceptance on 15 March 2024, and online publication on 11 March 2025. Its subscription PDF was not acquired. The author’s publication page links this title and DOI to the arXiv manuscript. The version, dates, DOI spelling, and frozen 2025 membership are retained; this report does not identify v1 with the journal final version. Every page location below refers to v1.\par
The intended setting is a compact complex manifold of dimension \(n>1\) with a nonzero parallel Lee form. The introductory dimension restriction is explicit. The distinction between \(c_1(X)=0\) in real de Rham cohomology and the sign of \(c_1^{\mathrm{BC}}(X)\) is essential. Quasi-regularity concerns the compactness of the one-dimensional complex Lee group, and orbifold canonical classes include the branch divisor. These distinctions are retained in the tests below.\par
\section*{Confirmed findings}
\subsection*{P1. The constructive direction of Theorem 2.4 does not produce the asserted metric (Major)}
Location. The last two paragraphs of the proof, p. 9: the positive line bundle, the norm-square potential, and the equality defining \(\Omega_F\).\par
First witness: positivity. Take \(\mathcal Y=\mathbb P^1\), \(\mathcal L=\mathcal O(1)\) with its positive Fubini–Study Hermitian metric, and \(P=(\mathcal L\smallsetminus0)/\langle\lambda\rangle\), \(0<|\lambda|<1\). This is a smooth compact principal elliptic bundle and is definite as required. In a standard local frame, its squared norm is \[f(z,w)=\frac{|w|^2}{1+|z|^2}.\] At \((z,w)=(0,1)\), the Hermitian matrix of \(f^{-1}i\partial\bar\partial f\) is \[\begin{pmatrix}-1&0\\0&1\end{pmatrix}.\] Thus the displayed form is indefinite. Interchanging the order of the two Dolbeault operators reverses both eigenvalues and does not restore positivity. The exact Hessian calculation was also executed with the existing symbolic runtime.\par
Second witness: descent, even after repairing positivity. Let \(E'=\mathbb C/(\mathbb Z+i\mathbb Z)\), let \(\tau\) be translation by \(1/2\), and let \(q:E'\to E'/\langle\tau\rangle\). Pull back an ample line bundle with a constant-curvature metric to \(\mathcal L\) on \(E'\). The pullback bundle and its elliptic quotient carry the order-two deck action. Start with its invariant normalized Vaisman form \[\Omega_0=\omega_0+\theta_0\wedge J\theta_0,\qquad
 \omega_0=-dJ\theta_0.\] For \(u(x+iy)=\cos(2\pi x)\) and sufficiently small \(\varepsilon\ne0\), perturb the positive line metric by \(e^{\varepsilon u}\) and use the reciprocal potential. Equivalently, the resulting normalized Vaisman forms have Lee forms \(\theta_\pm=\theta_0\pm\varepsilon\,du\) and positive transverse forms \(-dJ\theta_\pm\). The Lee vector is unchanged and remains Killing; \(\tau\) interchanges the two forms. Their arithmetic average is \[\frac{\Omega_++\Omega_-}{2}
   =\Omega_0+\varepsilon^2du\wedge Jdu ,
 \qquad
 -dJ\!\left(\frac{\theta_++\theta_-}{2}\right)
 +\frac{\theta_++\theta_-}{2}
       \wedge J\!\left(\frac{\theta_++\theta_-}{2}\right)
   =\Omega_0.\] The claimed equality therefore fails on a genuine allowed finite pullback cover. There is a stronger obstruction. Write \(\eta=\varepsilon^2du\wedge Jdu\) on the elliptic base. Since \(d\eta=0\), the unique Lee form of this average, if it is to satisfy \(d\Omega=\vartheta\wedge\Omega\), is \[\vartheta=\frac{\omega_0}{\omega_0+\eta}\,\theta_0.\] The coefficient is nonconstant, so \(d\vartheta\ne0\). The average is not even LCK.\par
Impact and counterevidence. These are failures of the actual metric construction and descent in a central equivalence used for the orbifold models and examples. This justifies Major severity for that proof component; it does not require the headline classification theorems to be false. Both witnesses satisfy the relevant positivity, compactness, and finite-action conditions. Conversely, they do not disprove the existence assertion of Theorem 2.4.\par
Repair boundary. For a positive line bundle use \(f=\|s\|_h^{-2}\), or use the norm-square potential on the negative dual line bundle, with consistent Lee normalization. For descent define a new form \[\widehat\Omega_F=-dJ\theta_F+\theta_F\wedge J\theta_F,
 \qquad \theta_F=|F|^{-1}\sum_{g\in F}g^*\theta.\] Its transverse part is the average of positive transverse forms. Its restriction in the Lee directions is unchanged, so it is positive, invariant, and has Killing Lee vector. It is therefore Vaisman. It must not be identified with the arithmetic average of the original metrics.\par
\subsection*{P2. The printed definition needs the nonzero-Lee restriction (Major)}
Location. The introductory definition, pp. 1–2; Definition 2.1, pp. 4–5; Theorems B and C, p. 2.\par
On \(X=(\mathbb C/(\mathbb Z+i\mathbb Z))^2\), for \(s>0\) put \[\Omega_s=\frac{i}{2}
   \left(s\,dz_1\wedge d\bar z_1+
                  s^{-1}dz_2\wedge d\bar z_2\right).\] These distinct flat metrics have \(\theta=0\), \(d\Omega_s=\theta\wedge\Omega_s\), \(\nabla\theta=0\), \(c_1^{\mathrm{BC}}=0\), the same Lee class, and the same volume. Thus the self-contained printed definition admits them, whereas Theorem B(1) asserts uniqueness. Moreover \(\operatorname{Aut}_0(X)=X\) has complex dimension two, contradicting the literal one-dimensional conclusion of Theorem C.\par
The strongest counterevidence is the paper’s explicit non-Kähler motivation and its repeated use of a one-dimensional Lee group. Those facts identify the intended strict setting. However, Definition 2.1 introduces normalization \(\|\theta\|=1\) as an additional condition; it does not remove zero-Lee metrics from the preceding definition. Under that explicit printed definition, the examples satisfy the actual hypotheses and contradict the stated conclusions of B(1) and C. The formal theorem-domain defect is therefore Major.\par
State \(\theta\ne0\) throughout, or explicitly adopt the strict non-Kähler convention. These examples do not refute the theorems with that intended additional restriction.\par
\subsection*{P3. The quotient theorem requires effective isotropy (Minor)}
Location. Effective chart convention, p. 5; Theorem 2.2, p. 7.\par
Let \(H=E=\mathbb C/(\mathbb Z+i\mathbb Z)\) and \(K=\langle1/2\rangle\). The action of \(H\) on \(X=(E/K)\times\mathbb P^1\) through its translation action on \(E/K\) is holomorphic, proper, and locally free, with stabilizer \(K\) at every point. Its quotient is \(\mathbb P^1\). An effective orbifold has trivial generic stabilizer; a principal \(H\)-orbibundle above its regular locus must have fiber \(H\). Here the fibers are \(E/K\), not \(H\). Indeed \(E/K\) has modulus \(2i\) after rescaling its lattice, whereas an \(\mathrm{SL}_2(\mathbb Z)\) transform of \(i\) has imaginary part \(1/(c^2+d^2)\) and cannot equal \(2\). The asserted principal structure over this quotient is therefore impossible.\par
Require the appropriate faithful slice action, or explicitly permit ineffective orbifolds and retain their generic isotropy. The effective Lee action used in the main applications does not have this counterexample: it is an actual subgroup of \(\operatorname{Aut}(X)\). The correction concerns the broader stated quotient theorem.\par
\subsection*{P4. Coarse and orbifold cohomology need a coefficient restriction (Minor)}
Location. The assertion following the definition of \(H^\bullet_{\mathrm{orb}}(\mathcal Y,R)\), p. 7.\par
Let \(\mathcal Y=[\mathbb P^1/\mu_2]\), where the nontrivial element acts by \(z\mapsto-z\), and use \(R=\mathbf F_2\). This is an effective compact complex orbifold with coarse space \(\mathbb P^1\). The classifying space is the Borel construction. A fixed point gives a section of its projection to \(B\mu_2\), so \(H^\bullet(B\mu_2,\mathbf F_2)\) injects into its cohomology. The periodic resolution has both differentials zero over \(\mathbf F_2\), hence this cohomology is nonzero in every nonnegative degree. The coarse sphere has no cohomology in degree three. Thus the displayed map cannot be an isomorphism for every field.\par
Characteristic zero suffices; more generally, invert all isotropy orders. The paper’s real and complex cohomological applications satisfy that condition. No counterexample to those applications is claimed.\par
\subsection*{P5. The anti-Lee vector is not automatically a lattice element (Minor)}
Location. The paragraph beginning with the integrality assertion for \(-JB\), p. 8.\par
For any strict quasi-regular Vaisman metric \(g\), the metric \(ag\), \(a>0\), has the same Lee form and Lee group but Lee vector \(B/a\). The lattice \(\Lambda\) of this group is fixed and discrete. The nonzero vectors \(-JB/a\) therefore cannot all lie in \(\Lambda\), although all these metrics meet the proof’s hypotheses. Concretely, on a scalar Hopf surface with \(g=g_{\rm Euclidean}/|z|^2\), one has \(B=-2r\partial_r\); the time-one flow of \(-JB\) is multiplication by \(e^{2i}\), not a real scalar deck power.\par
What is needed is a positive real multiple of a primitive lattice element, or an explicit period normalization. This retains the rational direction of the Chern class, the primitive generator, and the elliptic-group description. It corrects the unjustified metric-dependent integrality, not the existence of that direction.\par
\subsection*{P6. Flat principal monodromy has the wrong target group (Major)}
Location. The proof of the image criterion for \(\psi_*\), p. 9.\par
Principal \(G\)-bundle transition maps are translations in \(G\). They are not, in general, holomorphic group automorphisms of \(G\), which fix its identity. For an explicit topologically trivial but holomorphically nontrivial example, take the complex torus obtained from \(\mathbb C^2\) by the lattice generated by \[(1,0),\quad(i,\sqrt2),\quad(0,1),\quad(0,i)\] and project to \(E=\mathbb C/(\mathbb Z+i\mathbb Z)\). Its fiber is the same square elliptic curve \(G\), with translation monodromy \(1\mapsto0\), \(i\mapsto\sqrt2\pmod{\mathbb Z+i\mathbb Z}\). A holomorphic section would lift to \(w=az+b\), since its derivative is doubly periodic. Its period conditions require \(a\in\mathbb Z+i\mathbb Z\) and \(ia-\sqrt2\in\mathbb Z+i\mathbb Z\), which is impossible. Thus the bundle is not the trivial principal bundle. A transition that is both a group automorphism and commutes with every principal translation is the identity. The stated automorphism-valued construction cannot represent this example. Tensoring it with a positive line-induced elliptic bundle places it in the image-criterion argument with the same positive Chern class.\par
This is the inverse direction of the \(\psi_*\) image criterion used on the same page to construct the line-bundle quotient for \(q^*P\) in Theorem 2.4. The phrase as in the above discussion therefore depends on the invalid representation and lift. The defect is Major in that core construction, even though the image criterion itself has a direct repair: \(c^{\mathbb Z}(P_0)=0\) implies that \(P_0\) is the image of \(u\in H^1(\mathcal Y,\mathcal O_\mathcal Y)\) in the exponential exact sequence. The image of this same \(u\) in \(H^1(\mathcal Y,\mathcal O_\mathcal Y^*)\) maps to \(P_0\) under \(\psi_*\). This avoids both the wrong target and an unsupported lift of an arbitrary representation.\par
\subsection*{P7. The deformation tensor is not in the declared fixed bundle (Minor)}
Location. The opening projection formula in the proof of Theorem 3.1, p. 10; the graph identity on p. 11.\par
Use a constant real Beltrami parameter \(0<\mu<1\) on one factor of a complex two-torus. The deformed subspaces are \[T_t^{0,1}=\mathbb C(\bar\partial+\mu\partial),\qquad
 T_t^{1,0}=\mathbb C(\partial+\mu\bar\partial).\] The printed definition gives \[-\pi_t^{1,0}(\bar\partial)
   =\frac{\mu}{1-\mu^2}
                    (\partial+\mu\bar\partial).\] This vector has a nonzero old \((0,1)\) component. Consequently the declared tensor does not belong to \(A^{0,1}\otimes T^{1,0}\) for the fixed original structure, as required by the subsequent Maurer–Cartan and Hodge constructions. The example is an actual integrable deformation of a compact complex manifold.\par
Use instead the graph tensor obtained by inverting the old \((0,1)\) projection on \(T_t^{0,1}\) and then taking the old \((1,0)\) projection. In the example it is \(\bar\partial\mapsto\mu\partial\), and \(T_t^{0,1}=(1+\phi_t)T_0^{0,1}\). The invariant Kuranishi argument works with this standard tensor. This is a local coordinate/type correction; it does not refute the automorphism dimension theorem.\par
\subsection*{P8. Triviality of the entire compact subgroup action requires connectedness (Minor)}
Location. The trivial representation argument and equation (7), p. 10; Claim 1, p. 11.\par
The theorem allows an arbitrary compact complex subgroup \(H\subset\operatorname{Aut}_0(X)\). Take \(X=\mathbb P^1\times E\) and \(H=C_2\times E\), where \(C_2\) acts by \(z\mapsto-z\) and \(E\) acts by translations. This is a compact complex subgroup of positive dimension. The summand \[H^0(\mathbb P^1,T\mathbb P^1)\otimes H^1(E,\mathcal O_E)
       \ \subset\ H^1(X,TX)\] contains \(\partial_z\otimes\eta\), which changes sign under \(C_2\). It is not an invariant class. With an \(H\)-invariant metric, its harmonic representative also contradicts (7). A holomorphic character on a disconnected compact complex group need only be constant on each component.\par
Apply the argument to \(H^0\). Its dimension is the same as that of \(H\), so the desired dimension bound is retained. The final assertion with \(H=\operatorname{Aut}_0(X)\) already concerns a connected group. This is a false intermediate claim for the broader subgroup, not a counterexample to Theorem F.\par
\subsection*{P9. The negative-curvature proof has inconsistent signs and normalization (Minor)}
Location. The first paragraph of the proof of Theorem 6.2, p. 20; equation (22) and the contraction following (25), p. 21.\par
The proof chooses a positive form \(\eta\) representing \(-c_1^{\mathrm{BC}}(X)\) and then prescribes \(\rho_0=\eta\). But a Ricci form represents \(2\pi c_1^{\mathrm{BC}}(X)\), which has the opposite sign. This is not a permissible prescribed-Ricci target. For a concrete negative-type example, take a definite elliptic bundle over a curve \(C\) of genus two. Every transverse Ricci form has \[\int_C\rho=2\pi\deg TC=-4\pi<0,\] whereas a positive \(\eta\) has positive integral. All strict Vaisman and negative-class hypotheses are met.\par
Choose \(\eta\) in \(-2\pi c_1^{\mathrm{BC}}(X)\) and prescribe \(\rho_0=-\eta\). Equation (22) already uses this negative sign, which is strong counterevidence against treating the subsequent vanishing mechanism as false. A second displayed sign is also inconsistent: from the stated formula for \(\nabla^{\rm Ch}\theta^{1,0}\) and \(\theta^{1,0}(\xi_0)=0\), one obtains \[\theta^{1,0}(\nabla^{\rm Ch}\xi_0)
    =+\tfrac12 g(\xi_0,\cdot),\] rather than its negative. Only its squared norm is used.\par
With these corrections the independent computation gives \[\|D\xi_0\|_2^2=\|\xi_0\|_2^2-\|\xi_0\|_{h,2}^2,\qquad
 \|(\nabla^{\rm Ch}\xi_0)_0\|_2^2
                       =-\|\xi_0\|_{h,2}^2.\] Thus \(\xi_0=0\) and the intended one-dimensional automorphism conclusion is retained. Severity reflects the localized transcription and normalization defects, with the corrected core calculation actually checked.\par
\subsection*{P10. Continuity does not fix the Chern class of a chosen family of forms (Major)}
Location. The last cohomological inference in Corollary 6.6, p. 22.\par
On the constant family of curves \(C\times\Delta\) with \(g(C)>1\), choose a positive form \([\omega_0]=-c_1(C)\) and put \(\omega_t=(1+|t|^2)\omega_0\). This is a smooth family of fiberwise Kähler forms, obtained by restricting a smooth two-form on the total space, and has the prescribed central class. Nevertheless \([\omega_t]\ne-c_1(C)\) for \(t\ne0\). It satisfies exactly the properties of the chosen form stated before the claimed continuity inference.\par
This inference is the written step establishing negative type for the deformed fibers in Corollary 6.6, and thus the negative branch of Theorem E on p. 3. Fiberwise Kählerness and continuity do not supply the missing equality of cohomology classes; the source does not assume that the total form is closed or has constant Gauss–Manin class. This is a Major defect in the actual deformation-stability derivation.\par
The negative-type conclusion is recoverable: extend a positive Hermitian metric on the relative orbifold canonical bundle and use its fiberwise curvature. Its cohomology class is the required Chern class, and positivity persists after shrinking the parameter space. No counterexample to openness of negative Vaisman type is claimed.\par
\subsection*{P11. A properly elliptic Vaisman surface can have coarse base genus zero (Minor)}
Location. Example 7.1(1), p. 23.\par
Let \(\mathcal Y=\mathbb P^1(2,3,7)\). Its orbifold canonical degree is \[\deg K_\mathcal Y^{\mathrm{orb}}
   =-2+\left(1-\tfrac12\right)+\left(1-\tfrac13\right)
                               +\left(1-\tfrac17\right)
   =\frac1{42}>0.\] The local character of \(K_\mathcal Y^{\mathrm{orb}}\) is faithful at each orbifold point. Hence the nonzero total space of its negative dual is smooth. A negative-curvature Hermitian metric gives a positive norm-square cone metric, and quotienting by a real fiber contraction \(0<q<1\) gives a compact strict Vaisman surface \(X\). The circle connection makes \(\pi^*c_1(K_\mathcal Y^{\mathrm{orb}})\) exact, so \(c_1(X)=0\) in real cohomology. Since \(K_X=\pi^*K_\mathcal Y^{\mathrm{orb}}\) and the fibers are compact connected elliptic curves, pluricanonical sections descend; their growth gives \(\kappa(X)=1\). Thus \(X\) is properly elliptic but its coarse base is \(\mathbb P^1\), of genus zero.\par
The correct condition is positive orbifold canonical degree. Genus greater than one is available after a suitable finite orbifold uniformizing cover, not necessarily on the original coarse base. The negative Bott–Chern classification survives this correction.\par
\subsection*{P12. The punctured cone is the tautological bundle, not the degree-\(d\) bundle (Minor)}
Location. Example 7.4, p. 24, immediately before the adjunction argument.\par
Take weights \((1,1,1)\) and the smooth Fermat plane cubic \(f=z_0^3+z_1^3+z_2^3\), so \(d=3\). Its affine cone is smooth away from the origin, as required. The punctured cone is \(\mathcal O(-1)|_{Y_f}\) with its zero section removed. This line bundle has degree \(-3\). The printed \(\mathcal O(-d)|_{Y_f}=\mathcal O(-3)|_{Y_f}\) instead has degree \(-9\). Their circle bundles have abelianized fundamental groups \[\mathbb Z^2\oplus\mathbb Z/3\mathbb Z
       \quad\hbox{and}\quad
 \mathbb Z^2\oplus\mathbb Z/9\mathbb Z,\] respectively, so the claimed identification is genuinely false.\par
Replace \(\mathcal O(-d)\) by \(\mathcal O(-1)\) in the cone description. The corrected line bundle remains negative, and the adjunction condition \(d=\sum_iw_i\) is unaffected. The construction still yields the asserted Calabi–Yau examples. The problem is the specified bundle and its topology, not the existence of those examples.\par
\section*{Proof coverage and retained conclusions}
\subsection*{Orbifold correspondence and automorphisms}
The integral Chern class, branch divisor, and distinction between effective orbifolds and their coarse spaces were checked. The finite-cover part of Theorem 2.4 can be repaired without asserting that some finite cover has entirely torsion-free \(H_1\). For each \(m\)-torsion Chern component \(\alpha\), the Bockstein \[H^1_{\mathrm{orb}}(\mathcal Y,\mathbb Z/m\mathbb Z)\longrightarrow
                       H^2_{\mathrm{orb}}(\mathcal Y,\mathbb Z)[m]\] is onto. The cover defined by the kernel of the corresponding finite character kills \(\alpha\). Intersect the finitely many kernels for the lattice components. The exponential cohomology argument in P6 then gives the line-bundle model. The stronger torsion-free-\(H_1\) assertion was not declared false without a counterexample.\par
For connected compact \(H\), the holomorphic representations on \(H^q(X,TX)\) are trivial. Averaging the metric makes the harmonic projection and the Kuranishi construction equivariant. The invariant smooth auxiliary germ has the same tangent space as the full auxiliary germ, so the germs agree. The real \(H\) action persists under the corrected graph deformation, and the complex span of its Lie algebra has dimension at least \(\dim_\mathbb CH\). Upper semicontinuity gives equality when \(H=\operatorname{Aut}_0(X)\). P7–P8 isolate the necessary corrections; no counterexample to this argument with \(H^0\) was found.\par
\subsection*{Canonical metrics, signs, and deformation}
The Bott–Chern kernel description and the Lee cone explain the three signs; this use retains the strict Vaisman hypothesis. In Lemma 4.4, wedging the local transverse canonical form with \(\theta^{1,0}\) gives a holomorphic top form: the potentially nonzero \(\bar\partial\theta^{1,0}\) is horizontal and its wedge with the transverse top form vanishes. Its norm yields \(\rho_\Omega=\rho_0\). The positive-type Kodaira-dimension argument follows from Gauduchon integration of the divisor class against the Vaisman metric; the divisor integral need only be nonnegative.\par
The actual v2 of Ornea–Verbitsky [OV] was read and its Theorem 4.1 and proof pages were inspected. Its total-volume compatibility does not invalidate Theorem 5.1 here: first rescale the initial metric by a constant to match the target total volume. This changes neither the Lee class nor invariance under the Lee group. For \(c_1^{\mathrm{BC}}=0\), the Ricci potential is basic: its difference under an automorphism is pluriharmonic, hence constant, and boundedness under iteration makes that constant zero. The prescribed volume then produces the claimed Ricci-flat metric, and a pluriharmonic log ratio of two volumes proves uniqueness up to scale.\par
In the negative case, equations (12)–(18), the conformal adjoints, the \(nD_B\) cancellation, and the Chern orthogonal projection were checked directly. P9 records the two inconsistencies. The hypotheses of the transverse prescribed-Ricci theorem were checked on the retrieved pages of [E90]. In particular, the top basic class is nonzero: if \(\omega_0^{n-1}=d\beta\) were basic exact, integration against \(\theta\wedge J\theta\) and Stokes’ theorem would give zero, since \(\beta\wedge\omega_0\) exceeds the transverse dimension. This contradicts positivity of the total volume. The cited passage assumes transverse complex dimension at least two. In transverse dimension one, the required Ricci prescription follows directly by multiplying the transverse metric by the exponential of a basic potential and then normalizing its basic area. This boundary case was not silently attributed to the higher-dimensional statement.\par
Faulk’s actual Theorem 1.3 [Faulk] supplies the negative orbifold Kähler–Einstein metric for an effective compact Kähler orbifold. Uniqueness follows from the usual potential and maximum-principle argument. For Theorem 6.3, a closed \(G\)-invariant Lee candidate with the specified positive transverse curvature has positive coefficient in the old Lee direction; closedness eliminates the anti-Lee component, and pairing with a positive transverse power determines its positive sign. Thus the proposed form is positive and its Lee vector is Killing. The missing negative sign in P9, rather than a positive-type reading of this section, is the relevant correction. Theorem B explicitly imposes negativity.\par
\subsection*{Decomposition, topology, and examples}
The compactness of \(\operatorname{Aut}_0\) in the null case follows by pulling back the unique fixed-volume canonical metric. In the splitting argument, the parallel complex structure preserves the flat factor: a mixed component would give a nonzero parallel one-form on the simply connected compact factor. The primitive integral Chern class on that factor yields the cyclic root cover of the \(\mathbb C^*\)-bundle and the vanishing of its first homology. The retrieved Kollár preprint [Kollar], Proposition 88 on PDF p. 28, has the smooth-total-space, vanishing-\(H_1\), and numerical Calabi–Yau hypotheses needed to eliminate the branch divisor. Its numbering was not identified with the journal’s Proposition 10.2. The remaining pure-orbifold decomposition and unobstructedness inputs remain cited external theorems, not independently reproved results.\par
The finite canonical cover, the invariant part of a smooth Kuranishi germ, constancy of the holomorphic vector-field dimension, and the extension of a nonvanishing canonical form were checked as the dependency chain for Theorem 5.7. The small fibers are first kept Vaisman by the earlier deformation result. Their Frölicher degeneration and upper semicontinuity then give the required constancy of \(h^{n,0}\). The torsion canonical bundle gives Kodaira dimension zero; the elliptic fibration over a projective base and failure of the \(\partial\bar\partial\) lemma give algebraic dimension \(n-1\).\par
For a nontrivial simply connected factor, the fundamental group extension has central kernel \(\mathbb Z\oplus T\), \(T\) finite. Odd first Betti number forces the commutator to have rank zero, hence to lie in \(T\). If \(qT=0\), a subgroup generated by suitable \(q\)th powers of base lifts and a free central generator is abelian of finite index; passing to its normal core preserves abelianity. This supplies the needed finite-cover step without invoking an unsupported general virtual torsion-freeness claim. The first-Betti bound can also be obtained directly from the Chern-class sequence. Equality forces the finite base action on the abelian variety to be by translations, recovering a Kodaira manifold.\par
The weighted hypersurface examples were checked for quasi-smoothness, adjunction degree, and the displayed local stabilizer. Their gradients vanish only at the cone vertex. In the second threefold example the residual \(\mu_2\) acts by \((-1,-1)\) on the two local surface coordinates. P11 and P12 give the actual corrections to the base-genus and cone-bundle descriptions.\par
\section*{A conditional group-theoretic limitation}
Corollary 5.10, p. 18, requires special care if \(H_{\mathbb Z}\) means the usual unimodular integer Heisenberg group. On an abelian surface with polarization type \((1,2)\), the associated Kodaira threefold has group \(\Gamma=\mathbb Z_b\times\Gamma_B\), where \[B=\operatorname{diag}(J,2J),\qquad
 J=\begin{pmatrix}0&1\\-1&0\end{pmatrix},\] and the nontrivial basic commutators are \([x_1,y_1]=a\), \([x_2,y_2]=a^2\). This is a strict Vaisman Calabi–Yau example. Suppose a finite-index normal subgroup \(N\) were \(\mathbb Z\) times the standard unimodular Heisenberg group. Let \(V\subset\mathbb Z^4\) be its projected lattice. Then \(Z(N)=N\cap\langle a,b\rangle\). Its derived generator \(a^C\) is primitive in that full rank-two center, so \(N\cap\langle a\rangle=\langle a^C\rangle\). The form \(B/C\) must be unimodular on \(V\). Normality implies \[B(V,\mathbb Z^4)\subset C\mathbb Z,\qquad
 \mathbb Z^4\subset V^\#=V.\] Hence \(V=\mathbb Z^4\), but \(\det(B/C)=4/C^4\ne1\) for every positive integer \(C\). The extra central \(\mathbb Z\) therefore does not remove the obstruction. Dropping normality does: the subgroup generated by \(b,a^2,x_1^2,y_1,x_2,y_2\) is standard and of finite index, but is not normal.\par
This is a proved obstruction under the specified standard interpretation. It is not entered as an unqualified counterexample to the manuscript’s intended theorem. Example 7.2 itself calls a lattice integral Heisenberg when its commutator equals its center; the type-\((1,2)\) Heisenberg factor already satisfies that broader condition. No formal definition in v1 resolves this terminology. An exact integral classification must specify the lattice, allow the appropriate Heisenberg lattice in the normal subgroup, or assert standard commensurability without normality. The Betti conclusions do not require unimodularity. This interpretive boundary remains part of the pending scientific assessment.\par
\section*{Supporting checks}
All 1 477 normalized source lines and all 460 manifest lines were read by the initial root reviewer to zero remaining. The complete primary PDF was parsed and hash-bound. All used mathematical source pages were actually inspected by the root or the existing native bounded auxiliaries; their exact page/image bindings distinguish those roles. An index, extraction, or rendered image alone was not treated as scientific reading or visual inspection. Reviewer-authored exact calculations corroborate the Hessian, averaging identity, cyclic cohomology scalars, orbifold degree, polarization invariants, cone degrees, and displayed polynomial derivatives. No author code was run, and these small calculations are not represented as formal verification of the whole paper.\par
The initial root read their final notes and integrated the proofs, counterevidence, and limitations. Those initial checks did not constitute independent acceptance. The current focused acceptance and its native read-only contribution are separately attributed in verification records. No human review is claimed.\par
The directly retrieved external PDFs included Ornea–Verbitsky v2, Faulk v1, El Kacimi-Alaoui’s original article, and the Kollár author preprint. Their full programmatic parsing, bounded scientific reading, and actual image inspections are recorded separately. The Borzellino–Zhu PDF retrieval attempts failed through the documented routes; indexed theorem text was not relabeled as an inspected PDF. Campana’s decomposition theorem, the cited unobstructedness and Frölicher results, and the complete proofs of other classical inputs were not all independently reproved. These are substantive dependency limits, not confirmed mathematical defects.\par
The publisher, arXiv, and author’s publication page were searched for this exact article and for supplementary material, code, and corrections [Istrati, EMS, Author]. No separate paper-specific supplement, implementation, or correction was found in that bounded search. Generic code-provider links are not associated code. The absence of code is not a defect of this mathematical paper. No subscription access was bypassed, and the uninspected journal version is outside the mathematical claims of this report.\par
The current TeX, actual compiler execution, dependency snapshot, PDF, and all final report-page inspections are bound in the delivery records. Earlier failures, drafts, raw inspection-false facts, and withdrawn candidate assessments remain historical records. A fresh current record controls delivery, without rewriting those facts.\par
The current assessment is four Major and eight Minor findings. The scientific decision remains pending, with `human\_reviews=[]`; the independent acceptance is recorded separately against the exact current artifacts.\par
9 N. Istrati, Vaisman manifolds with vanishing first Chern class, arXiv:2304.02582v1 (2023), frozen author source, https://arxiv.org/abs/2304.02582v1. L. Ornea and M. Verbitsky, A Calabi–Yau theorem for Vaisman manifolds, arXiv:2206.08808v2 (2022), Theorem 4.1, PDF pp. 7–9, https://arxiv.org/abs/2206.08808v2. A. El Kacimi-Alaoui, Opérateurs transversalement elliptiques sur un feuilletage riemannien et applications, Compositio Math. 73 (1990), 57–106, §3.5.5; original scanned article. M. Faulk, On Yau’s theorem for effective orbifolds, arXiv:1710.11561v1 (2017), Theorem 1.3, PDF p. 2, https://arxiv.org/abs/1710.11561v1. J. Kollár, Einstein metrics on 5-dimensional Seifert bundles, arXiv:math/0408184v1, Proposition 88, PDF p. 28, https://arxiv.org/abs/math/0408184v1. EMS Press, publisher article record for DOI 10.4171/JEMS/1633, accessed 27 October 2026, https://ems.press/journals/jems/articles/14298650. N. Istrati, research publication page, accessed 27 October 2026, https://sites.google.com/site/nicolinaistrati/research.\par
\end{document}
